\zadanie{Układ o równoodległych poziomach energetycznych}

Rozważamy opisany rozkładem kanonicznym układ znajdujący się w równowadze ze
zbiornikiem ciepła o temperaturze $T$. Poziomy energetyczne układu dane są przez
\begin{equation*}
  \varepsilon_j = j\,\varepsilon, \qquad j = 0, 1, 2, 3, \ldots, \infty ,
\end{equation*}
gdzie $\varepsilon > 0$ jest odległością między sąsiednimi poziomami
energetycznymi. Wiemy, że średnia energia tego układu to
\begin{equation*}
  \langle E \rangle = E(T) = \alpha\varepsilon ,
\end{equation*}
gdzie $\alpha > 0$ jest pewnym wyrażeniem.

\begin{enumerate}[label=(\alph*)]
  \item Pokaż, że temperatura $T$ rozważanego układu spełnia równanie
    \begin{equation*}
      \frac{1}{k_B T} = \frac{1}{\varepsilon}\ln\left(1 + \frac{1}{\alpha}\right).
    \end{equation*}
  \item Pokaż, że entropia układu dana jest wyrażeniem
    \begin{equation*}
      S = k_B(1 + \alpha)\ln(1 + \alpha) - k_B\,\alpha\ln\alpha .
    \end{equation*}
\end{enumerate}

\begin{wskazowki}
\begin{enumerate}[label=(\alph*)]
  \item \textit{Temperatura.}
    \begin{itemize}
      \item Zapisz kanoniczną funkcję podziału układu jako sumę po poziomach
        energetycznych:
        \begin{equation*}
          Z = \sum_{j=0}^{\infty} e^{-\varepsilon_j/k_BT} .
        \end{equation*}
      \item Kolejne wyrazy tej sumy tworzą szereg geometryczny. Jaki jest jego
        iloraz? Skorzystaj ze wzoru
        $\sum_{j=0}^{\infty} q^j = 1/(1 - q)$ dla $|q| < 1$.
      \item Korzystając z $Z$, oblicz średnią energię:
        \begin{equation*}
          E(T,V,N) = k_BT^2\,\frac{\partial}{\partial T}\ln Z(T,V,N) .
        \end{equation*}
    \end{itemize}
  \item \textit{Entropia.}
    \begin{itemize}
      \item Skorzystaj ze wzoru $S = k_B\ln Z + E/T$.
      \item Dzięki wynikowi punktu~(a) możesz wyrazić $\ln Z$ i $E/T$
        przez $\alpha$ zamiast przez $T$.
    \end{itemize}
\end{enumerate}
\end{wskazowki}
